% =============================================================================
% GLOSSAIRE DES SIGLES
% =============================================================================
\chapter*{Glossaire des sigles}
\addcontentsline{toc}{chapter}{Glossaire des sigles}

\begin{description}[leftmargin=3.5cm,style=nextline]
  \item[ACO] \emph{Ant Colony Optimization} — algorithme d'optimisation par colonies de fourmis (Dorigo, 1992).
  \item[ADR] \emph{Architectural Decision Record} — document tra\c{c}ant une d\'ecision d'architecture et son rationnel.
  \item[Agentic BPM] Gestion de processus m\'etier o\`u les acteurs peuvent \^etre des agents logiciels autonomes.
  \item[BPMN] \emph{Business Process Model and Notation}.
  \item[CAS] \emph{Complex Adaptive Systems} (Holland, 1995).
  \item[DSR] \emph{Design Science Research} (Hevner et~al., 2004 ; Peffers et~al., 2007).
  \item[EU AI Act] R\`eglement europ\'een sur l'intelligence artificielle, adopt\'e en 2024.
  \item[FEDS] \emph{Framework for Evaluation in Design Science} (Venable et~al., 2016).
  \item[JSONL] \emph{JSON Lines} — format texte append-only, une ligne JSON par enregistrement.
  \item[LLM] \emph{Large Language Model} — grand mod\`ele de langage.
  \item[MAS] \emph{Multi-Agent System}.
  \item[MAST] \emph{Multi-Agent System Failure Taxonomy} (Cemri, 2025) — 14 modes d'\'echec class\'es.
  \item[MHC] \emph{Meaningful Human Control} — supervision humaine effective d'un syst\`eme automatis\'e.
  \item[MIS] \emph{Management Information Systems}.
  \item[OC] Objectif de Conception (DO en anglais : \emph{Design Objective}).
  \item[RPA] \emph{Robotic Process Automation}.
  \item[SQLite] Moteur de base de donn\'ees relationnel embarqu\'e en mode WAL utilis\'e pour persister les markers.
  \item[TTL] \emph{Time-To-Live} — dur\'ee de vie maximale d'un verrou ou d'un marker actif.
  \item[WAL] \emph{Write-Ahead Logging} — mode de journalisation SQLite garantissant durabilit\'e et concurrence.
\end{description}

\cleardoublepage
